\setcounter{page}{11}
\chapter{Operaciones vectoriales}
\label{cap:cap2}
\vspace{-15ex}

\raggedright


En SEMANA \ref{cap:cap1} mencionamos que aquellas cantidades f\'isicas que requieren no solo de una magnitud, sino de una direcci\'on para especificarse, se llaman \textit{cantidades vectoriales}, o simplemente \textit{vectores}. En esta sesi\'on estudiaremos la manera en la cual se expresan y representan los vectores; así como la manera de realizar operaciones matemáticas con ellos.

\section{Notaci\'on vectorial y vectores unitarios}

\begin{comment}
Si observamos, por ejemplo, una imagen satelital del campus central del TEC (Figura \ref{fig:campus}), podemos decir que \underline{respecto} al \textsf{Cajero BCR}, el \textsf{Centro de las Artes} se encuentra aproximadamente hacia el noreste (NE). Utilizando la escala que aparece en el margen inferior derecho de la imagen, podemos agregar que el \textsf{Centro de las Artes} se encuentra a una distancia de aproximadamente 275 m del \textsf{Cajero BCR}.\\[.35cm]
\begin{figure}
\input{ejemplo_vector_posicion}
\caption{Campus Central del TEC, Cartago.}
\label{fig:campus}
\end{figure}
\end{comment}

%Por lo tanto, la posici\'on de un objeto (respecto a alg\'un punto de referencia) es un vector. Llamemos $\vec{r}$ al vector que representa la posici\'on. 
%Para el caso ejemplificado $$\vec{r}_{\Mbox{\textsc{Centro de las Artes}}}=(275\U{m};45^{\circ}), \quad \mbox{respecto al \textsc{Cajero BCR.}}$$

%Se le coloca una flecha ($\vec{\phantom{a}}$) encima del s\'imbolo que representa una cantidad f\'isica vectorial para indicar su naturaleza direccional.\\[.35cm]
 
Podemos representar un vector haciendo uso de las dos propiedades b\'asicas de un vector: su magnitud y su direcci\'on. Si ``A'' es una cantidad vectorial, de magnitud $A$ y direcci\'on $\theta$, respecto a un sistema de referencia, entonces 
$$\vec{A}=(A;\theta) \qquad \mbox{(\textsf{notaci\'on polar de un vector)}}. $$

En la Figura \ref{fig:polar} se muestra un \textit{sistema de referencia polar}.\\[.35cm]

\begin{figure}[!ht]
%\begin{center}
%\begin{tabular}[cc]
\input{polar}  % \input{polarplane}
%\end{tabular}
\caption{Sistema de coordenadas polar.}
\label{fig:polar}
%\end{center}
\end{figure}


Otra manera de representar un vector, es indicando sus \textit{componentes cartesianas} o \textit{rectangulares} respecto a un \textit{sistema de referencia cartesiano}. \\[.35cm]

 Un \textit{sistema cartesiano} es un conjunto de ejes mutuamente perpendiculares. En la Figura \ref{fig:sist} se muestra un \textit{sistema de coordenadas cartesiano} y el \textit{vector de posici\'on} de un punto con coordenadas ($x_0,y_0,z_0$). Las cantidades $\hat{\iota}$, $\hat{\jmath}$ y $\hat{\kappa}$ se llaman \textit{vectores directores unitarios} y son vectores de magnitud uno y direcci\'on a lo largo de los ejes $x$, $y$ y $z$, respectivamente.  \\[.35cm]
 
 \begin{figure}[!ht]
\begin{center}
\input{marco_referencia}
\caption{Sistema de coordenadas cartesiano. $\vec{r}$  es el \textit{vector posici\'on} del punto ($x_0,y_0,z_0$).}
\label{fig:sist}
\end{center}
\end{figure}

 
 Por lo tanto, para el caso del vector de posici\'on de un punto con coordenadas cartesiandas $(x_{0},y_{0},z_{0})$ se denota 
$$\vec{r}=x_{0}\hat{\iota}+y_{0}\hat{\jmath}+z_{0}\hat{\kappa} \qquad \mbox{\textsf{(notaci\'on cartesiana, vector de posici\'on).}} $$



Dado un punto en un plano (espacio 2D) con coordenadas cartesianas $(x,y)$, es posible obtener las correspondientes coordenadas polares mediante la transformaci\'on
$$ r=\sqrt{x^2+y^2} \qquad \mbox{y} \qquad \theta= \arctan \left( \frac{y}{x}\right), $$

o de forma inversa

$$x=r \cos \theta \qquad \mbox{y} \qquad y=r \sin \theta.$$

Los puntos cardinales son las cuatro direcciones principales que se utilizan para la navegación y la orientación: Norte (\textbf{N}), Sur (\textbf{S}), Este (\textbf{E}) y Oeste (\textbf{O}):

\begin{itemize}
    \item Norte (N): Es la dirección hacia el Polo Norte. En los mapas, el norte suele estar en la parte superior.
    \item Sur (S): Es la dirección hacia el Polo Sur, opuesta al norte.
    \item Este (E): Es la dirección por donde sale el sol. En los mapas, el este está a la derecha.
    \item Oeste (O): Es la dirección por donde se pone el sol, opuesta al este. En los mapas, el oeste está a la izquierda.
\end{itemize}
\begin{comment}
Para el ejemplo del $\vec{r}_{\Mbox{\textsc{Centro de las Artes}}}$ ver Figura \ref{fig:ca}.

\begin{figure}[ht!]
\centering
\begin{tikzpicture}
\input{ejemplo_componentes}
\end{tikzpicture}
\caption{Relaci\'on entre notaci\'on polar y cartesiana.}
\label{fig:ca}
\end{figure}
\end{comment}

\begin{figure}
    \centering
    \includegraphics[width=5in]{planeamiento_didactico/figuras/rosa_vientos.pdf}
    \caption{Rosa de los vientos.}
    \label{fig:rosa-vientos}
\end{figure}

\begin{figure}
    \centering
    \includegraphics[width=5in]{planeamiento_didactico/figuras/puntos_cardinales.pdf}
    \caption{Direcci\'on cardinal vs direcci\'on polar. Para ubicaciones en el I y II cuadrantes, se toma el N como referencia.}
    \label{fig:puntos-cardinales}
\end{figure}

\begin{figure}
    \centering
    \includegraphics[width=5in]{planeamiento_didactico/figuras/puntos_cardinales-2.pdf}
    \caption{Direcci\'on cardinal vs direcci\'on polar.  Para ubicaciones en el III y IV cuadrantes, se toma el S como referencia.}
    \label{fig:puntos-cardinales-2}
\end{figure}
\section{Operaciones vectoriales}

Partiendo de la notaci\'on cartesiana
\begin{align*}
\vec{A}=&(A_x,A_y,A_z)\\[.35cm]
=&A_x\hat{\iota}+A_y\hat{\jmath}+A_z\hat{\kappa}.
\end{align*}
definimos las siguientes operaciones vectoriales:
\resp{.75\textwidth}{\textbf{Producto de un escalar y un vector:}
\begin{align*}
a\vec{A}=&a(A_x,A_y,A_z)\\[.35cm]
=&(aA_x,aA_y,aA_z) \\[.35cm]
=&aA_x\hat{\iota}+aA_y\hat{\jmath}+aA_z\hat{\kappa},
\end{align*}}{white}{noticias}

\begin{wwteorema}{Ejemplo}
La fuerza neta ($\vec{F}_{\Mbox{neta}}$) que experimenta un cuerpo de masa $m$ se relaciona con la aceleraci\'on ($\vec{a}$): $$\vec{F}_{\Mbox{neta}}=m\vec{a}$$
Si sobre un cuerpo de $m=$10 kg act\'ua una fuerza neta de $$\vec{F}_{\Mbox{neta}}=(25\U{N})\hat{\iota}-(75\U{N})\hat{\jmath}+(100\U{N})\hat{\kappa},$$
este experimentar\'a una aceleraci\'on de 
$$\vec{a}=\left(\frac{1}{m}\right)\vec{F}_{\Mbox{neta}}=(2.5\U{m/s}^2)\hat{\iota}-(7.5\U{m/s}^2)\hat{\jmath}+(10\U{m/s}^2)\hat{\kappa}$$
\end{wwteorema}


\resp{.75\textwidth}{
%\begin{itemize}
%\item 
\textbf{Suma (resta) de vectores}, $\vec{A} \pm \vec{B}:$
\begin{align*}
\vec{A} \pm \vec{B}=&(A_x,A_y,A_z)\pm(B_x,B_y,B_z)\\[.35cm]
=&(A_x \pm B_x,A_y \pm B_y,A_z \pm B_z) \\[.35cm]
=&(A_x \pm B_x)\hat{\iota}+(A_y \pm B_y)\hat{\jmath}+(A_z \pm B_z)\hat{\kappa}.
\end{align*}}{white}{noticias}

\begin{wwteorema}{Ejemplo}
La fuerza neta ($\vec{F}_{\Mbox{neta}}$) que experimenta un cuerpo est\'a dada por 
$$\vec{F}_{\Mbox{neta}}=\sum \vec{F}_i=\vec{F}_1+\vec{F}_2+\cdots + \vec{F}_N$$
Si sobre un objeto act\'uan dos fuerzas:
\begin{align*}
\vec{F}_1&=&(5\U{N})&\hat{\iota}&-&(15\U{N})\hat{\jmath}&+&(10\U{N})\hat{\kappa} \\[.35cm]
\vec{F}_2&=&(-15\U{N})&\hat{\iota}&+&(25\U{N})\hat{\jmath}&+&(10\U{N})\hat{\kappa} \\[.35cm]
\Rightarrow \vec{F}_{\Mbox{neta}}&=&(-10\U{N})&\hat{\iota}&+&(10\U{N})\hat{\jmath}&+&(20\U{N})\hat{\kappa} \\[.35cm]
\end{align*}
\end{wwteorema}

\resp{.75\textwidth}{
%\begin{itemize}
%\item 
\textbf{Producto escalar de dos vectores}, $\vec{A}\cdot\vec{B}:$
\begin{align*}
\vec{A}\cdot\vec{B}=&A_xB_x+A_yB_y+A_zB_z\\[.35cm]
=&|\vec{A}||\vec{B}|\cos \theta\\[.35cm]
=&\vec{B}\cdot\vec{A}.
\end{align*}
}{white}{noticias} 

donde $|\vec{A}|$ y $|\vec{B}|$ representan las magnitudes de $\vec{A}$ y $\vec{B}$ respectivamente, y $\theta$ el \'angulo entre ellos (ver Figura \ref{fig:apb}). Tambi\'en puede escribirse $|\vec{A}|=A$.\\[.35cm]

El producto escalar entre dos vectores, tambi\'en llamado \textbf{producto punto} es una cantidad escalar. \\[.35cm]
\begin{figure}[h!]
\begin{center}
\input{angulo_entre_vectores}
\end{center}
\caption{\'Angulo entre dos vectores.}
\label{fig:apb}
\end{figure}

\begin{wwteorema}{Ejemplo}
El trabajo ($W$) de un fuerza ($\vec{F}$) cuando act\'ua a lo largo de un desplazamiento ($\vec{s}$) est\'a dado por 
$$W=\vec{F}\cdot \vec{s}=Fs\cos \theta,$$
donde $\theta$ es el \'angulo entre $\vec{F}$ y $\vec{s}$.
\end{wwteorema}

Note que 
$$|\vec{A}|=\sqrt{\vec{A}\cdot\vec{A}}= \sqrt{A_x^{2}+A_y^{2}+A_z^{2}}.$$
Adem\'as, si dos vectores son perpendiculares, es decir si $\theta=90^{\circ}$, $\vec{A}\cdot\vec{B}=0$. \\[.35cm] 
A partir de esta operaci\'on, se pueden calcular los \textit{\textcolor{c5}{cosenos directores}} de un vector
\begin{align*}
\vec{A}\cdot\hat{\iota}=&|\vec{A}||\hat{\iota}|\textcolor{c5}{\cos\theta_{x}}=A\textcolor{c5}{\cos\theta_{x}}=A_{x}, \\[.35cm]
\vec{A}\cdot\hat{\jmath}=&|\vec{A}||\hat{\jmath}|\textcolor{c5}{\cos\theta_{y}}=A\textcolor{c5}{\cos\theta_{y}}=A_{y}, \\[.35cm]
\vec{A}\cdot\hat{\kappa}=&|\vec{A}||\hat{\kappa}|\textcolor{c5}{\cos\theta_{z}}=A\textcolor{c5}{\cos\theta_{z}}=A_{z},
\end{align*}
donde $\theta_{x}$, $\theta_{y}$ y $\theta_{z}$, llamados \textit{\'angulos directores}, son los \'angulos que forma el vector  con los ejes $x$, $y$ y $z$ respectivamente.





\resp{.75\textwidth}{\textbf{Producto vectorial de dos vectores}, $\vec{A}\times\vec{B}=\vec{C}:$
\begin{align*}
C_x=&A_yB_z-A_zB_y,\\[.35cm]
C_y=&A_zB_x-A_xB_z,\\[.35cm]
C_z=&A_xB_y-A_yB_z.
\end{align*}}{white}{noticias}

%\item 

El producto vectorial (tambi\'en llamado \textbf{producto cruz}) entre dos vectores produce un vector que es perpendicular al plano que contiene a ambos vectores (ver Figura \ref{fig:axb}). \\[.35cm] 
Se cumple adem\'as que 
$$ |\vec{A}\times\vec{B}|=|\vec{A}||\vec{B}|\sin \theta,$$
siendo $\theta$ el \'angulo entre $\vec{A}$ y $\vec{B}$.  Geom\'etricamente,  esta magnitud equivale al \'area del paralelogramo que se forma entre $\vec{A}$ y $\vec{B}$ (ver Figura \ref{fig:axb}). \\[.35cm] 

\input{producto_cruz}

Note que el producto cruz entre vectores paralelos $(\theta=0)$ es nulo. \\[.35cm] 

\begin{wwteorema}{Ejemplo}
La torca ($\vec{\tau}$) de un fuerza ($\vec{F}$) cuando act\'ua a una distancia ($\vec{r}$) del eje de giro de un cuerpo r\'igido est[a dado por 
$$\vec{\tau}=\vec{F}\times \vec{r}$$ \\[.35cm]
Si el eje de giro ($z$) permanece constante, 
$$\tau_z=Fr\sin \theta=Fr_{\perp},$$
donde $r_{\perp}$ se conoce como \textit{brazo de palanca} de la fuerza.
\end{wwteorema}
